\section{更深一点的局限：为什么榜单会失真}

这一讲反复提醒“evaluation crisis”，原因就是榜单会慢慢失真。

\subsection{污染与饱和}

一旦训练数据和测试数据有重叠，benchmark 就不再纯净。另一方面，数据集越简单，模型越快饱和，榜单越不能区分模型。

\subsection{奖励错配}

如果 leaderboard 奖励的是某种局部优势，比如固定格式、特定题型或短期偏好，模型开发就会围着这些奖励转，最后不一定更有用。

\subsection{judge bias}

LLM-as-judge 或 human judge 都可能偏向某些风格。比如更长、更流畅、更像“礼貌助手”的答案，未必真的更正确。

\begin{warningbox}{榜单不能替代理解}
榜单只能告诉你“这个系统在某些定义下表现更好”，不能自动告诉你“它为什么好”或“它在你关心的真实任务里也一定更好”。
\end{warningbox}

